% --- Chapter appendices: numbered 2.A, 2.B ---------------------------------
\renewcommand{\thesection}{\thechapter.\Alph{section}}
\setcounter{section}{0}

\section{An integral identity for the hypergeometric function}
\label{m:app:hyp}

The route to \cref{m:eq:Psidef} taken in \cref{m:sec:absbirthdeath} expands the
characteristic $x(\tau)$ in a geometric series and integrates term by term. An
alternative is to integrate the quotient \cref{m:eq:xabd} directly, splitting the
numerator into two pieces, each of the form
\begin{equation}
\label{m:eq:Iform}
I(\tau) = \int\frac{\ee^{h\tau}}{P+Q\,\ee^{k\tau}}\,\dd\tau
\end{equation}
with $h,k,P,Q$ constant. That route is longer and easier to get wrong --- the two
pieces carry different prefactors, and dropping one produces a formula that still
satisfies the initial condition while violating $G(1,1,t)=1$ --- but the identity it
rests on is useful in its own right, so it is recorded and proved here.

\begin{proposition}
\label{m:prop:hypIden}
For $h/k>0$ and $-\tfrac{Q}{P}\ee^{k\tau}\notin[1,\infty)$,
\begin{equation}
\label{m:eq:hypIden}
\int\frac{\ee^{h\tau}}{P+Q\ee^{k\tau}}\,\dd\tau
= \frac{\ee^{h\tau}}{Ph}\,{}_2F_1\lrb{1,\frac{h}{k};\frac{h}{k}+1;\,-\frac{Q}{P}\ee^{k\tau}}
+ \text{const} .
\end{equation}
\end{proposition}

\begin{proof}
We use the Euler integral representation of the hypergeometric function
\cite{Olver2010},
\begin{equation}
\label{m:eq:idenHyp}
{}_2F_1(a,b;c;z)
= \frac{\Gamma(c)}{\Gamma(b)\Gamma(c-b)}\int_0^1 u^{b-1}(1-u)^{c-b-1}(1-zu)^{-a}\,\dd u ,
\end{equation}
$\Gamma$ being the gamma function.

Start from \cref{m:eq:Iform} and substitute $v = \ee^{k\tau}$, so that
$\dd\tau = k^{-1}v^{-1}\dd v$:
\begin{equation*}
I = \frac1k\int\frac{v^{h/k}}{P+Qv}\,v^{-1}\dd v
= \frac{1}{Pk}\int\frac{v^{h/k-1}}{1+\frac{Q}{P}v}\,\dd v .
\end{equation*}
Fix the antiderivative by writing this as a definite integral from $0$ to $v$ in a
dummy variable $w$; the integral converges at the lower limit precisely because
$h/k>0$:
\begin{equation*}
I = \frac{1}{Pk}\int_0^v\frac{w^{h/k-1}}{1+\frac{Q}{P}w}\,\dd w .
\end{equation*}
Now rescale, $w = vu$, so that $\dd w = v\,\dd u$ with $v$ held fixed:
\begin{equation}
\label{m:eq:backat}
I = \frac{1}{Pk}\int_0^1\frac{(vu)^{h/k-1}}{1+\frac{Q}{P}vu}\,v\,\dd u
= \frac{v^{h/k}}{Pk}\int_0^1\frac{u^{h/k-1}}{1+\frac{Q}{P}vu}\,\dd u .
\end{equation}
The remaining integrand matches that of \cref{m:eq:idenHyp} with $b = h/k$,
$c = h/k+1$, $a=1$ and $z = -\tfrac{Q}{P}v$, since then $(1-u)^{c-b-1} = 1$ and
$(1-zu)^{-a} = \bigl(1+\tfrac{Q}{P}vu\bigr)^{-1}$. Hence, using
$\Gamma(z+1) = z\Gamma(z)$,
\begin{equation*}
\int_0^1\frac{u^{h/k-1}}{1+\frac{Q}{P}vu}\,\dd u
= \frac{\Gamma\!\lrb{\frac hk}\Gamma(1)}{\Gamma\!\lrb{\frac hk+1}}\,
{}_2F_1\lrb{1,\frac hk;\frac hk+1;-\frac QP v}
= \frac kh\,{}_2F_1\lrb{1,\frac hk;\frac hk+1;-\frac QP v}.
\end{equation*}
Substituting into \cref{m:eq:backat} gives
\begin{equation*}
I = \frac{v^{h/k}}{Ph}\,{}_2F_1\lrb{1,\frac hk;\frac hk+1;-\frac QP v},
\end{equation*}
and restoring $v = \ee^{k\tau}$ gives \cref{m:eq:hypIden}.
\end{proof}

\begin{remark}
Applied to the two pieces of \cref{m:eq:xabd} --- the first with $h=\alpha$, the
second with $h = \alpha+b_1$, and both with $k=b_1$, $P = \sigma-x_-$ and
$Q = -(\sigma-x_+)$ --- \cref{m:prop:hypIden} reproduces \cref{m:eq:Psihyp} once the
prefactors $\alpha x_+(\sigma-x_-)$ and $-\alpha x_-(\sigma-x_+)$ are carried
through correctly. The factor $(\sigma-x_+)/(\sigma-x_-)$ attached to the second
piece is easily lost, and its loss is invisible at $t=0$; checking $G(1,1,t)=1$
catches it immediately.
\end{remark}

\section{Extracting state probabilities from a generating function}
\label{m:app:extract}

A generating function defined by a formula rather than as an explicit series --- for
example
\begin{equation}
\label{m:eq:examplePgf}
G(x,y,t) = \Bigl(\bigl(1-\ee^{-\alpha t}\bigr)x + \ee^{-\alpha t}y\Bigr)^n
\end{equation}
--- can always be converted back to the form
\begin{equation}
\label{m:eq:series1}
G(x,y,t) = \sum_{i=0}^{\infty}\sum_{j=0}^{\infty}p_{i,j}(t)\,x^iy^j ,
\end{equation}
recovering the time-dependent state probabilities. All that is required is that $G$
be repeatedly partially differentiable in $x$ and $y$.

The reason is the bivariate Maclaurin expansion,
\begin{equation*}
f(x,y) = \sum_{k=0}^{\infty}\frac{1}{k!}\sum_{\varrho=0}^{k}\binom{k}{\varrho}
\partial_x^{\varrho}\partial_y^{\,k-\varrho}f\Big|_{(0,0)}\;x^{\varrho}y^{k-\varrho},
\end{equation*}
whose coefficient of $x^iy^j$, taking $k=i+j$ and $\varrho=i$, is
$\frac{1}{(i+j)!}\binom{i+j}{i}\partial_x^i\partial_y^{\,j}f\big|_{(0,0)}$. Hence
\begin{equation}
\label{m:eq:stateProb}
p_{i,j}(t) = \frac{1}{i!\,j!}\,\partial_x^i\partial_y^{\,j}G\Big|_{(0,0)} .
\end{equation}

\subsection*{The multilinear case}

Both \cref{m:eq:Gabsonly} and \cref{m:eq:Gabsdeath} have the form
\begin{equation}
\label{m:eq:multilinear}
G(x,y,t) = \bigl(c(t) + f(t)\,x + g(t)\,y\bigr)^n ,
\end{equation}
for which the derivatives are immediate. Each differentiation lowers the power by
one and brings down a factor, so
\begin{equation}
\label{m:eq:derivs}
\partial_x^{a}\partial_y^{\,b}G
= \frac{n!}{(n-a-b)!}\,f^{a}g^{b}\bigl(c+fx+gy\bigr)^{n-a-b}
\qquad(a+b\le n),
\end{equation}
and zero for $a+b>n$. Evaluating at the origin and applying \cref{m:eq:stateProb},
\begin{equation}
\label{m:eq:stateProbSpec}
p_{i,j}(t) = \frac{n!}{i!\,j!\,(n-i-j)!}\;f(t)^{i}g(t)^{j}c(t)^{\,n-i-j},
\qquad i+j\le n,
\end{equation}
and $p_{i,j}=0$ otherwise. This is the trinomial distribution, which is exactly
what it should be: the $n$ particles are independent, and at time $t$ each is in
one of three states, with probabilities $c$ (neither inside nor outside --- that
is, dead), $f$ (inside) and $g$ (outside). Note in particular that the falling
factorial $n!/(n-i-j)!$ in \cref{m:eq:derivs} is \emph{not} $n^{i+j}$; the two agree
only to leading order in $n$, and using the latter breaks normalisation.

\subsection*{Applied to the absorption models}

For absorption only, \cref{m:eq:examplePgf}, we have $c=0$, $f = 1-\ee^{-\alpha t}$
and $g = \ee^{-\alpha t}$. The factor $c^{\,n-i-j}$ then vanishes unless $i+j=n$,
and \cref{m:eq:stateProbSpec} gives
\begin{equation*}
p_{i,j}(t) = \binom{n}{i}\bigl(1-\ee^{-\alpha t}\bigr)^{i}\bigl(\ee^{-\alpha t}\bigr)^{j},
\qquad i+j=n,
\end{equation*}
so $X_t\sim\text{Binomial}\bigl(n,1-\ee^{-\alpha t}\bigr)$ and $Y_t = n-X_t$, as
anticipated in \cref{m:sec:absonly}.

For absorption--death, \cref{m:eq:Gabsdeath}, the three functions are
\begin{equation*}
c(t) = 1-\frac{\mu}{\mu-\alpha}\ee^{-\alpha t} + \frac{\alpha}{\mu-\alpha}\ee^{-\mu t},
\qquad
f(t) = \frac{\alpha}{\mu-\alpha}\bigl(\ee^{-\alpha t}-\ee^{-\mu t}\bigr),
\qquad
g(t) = \ee^{-\alpha t},
\end{equation*}
which are precisely $p_{0,0}$, $p_{1,0}$ and $p_{0,1}$ of
\cref{m:eq:absdeathstates}, so that
\begin{equation*}
p_{i,j}(t) = \frac{n!}{i!\,j!\,(n-i-j)!}\;f(t)^{i}g(t)^{j}c(t)^{\,n-i-j} .
\end{equation*}
Each particle independently ends up dead, absorbed, or still outside, with those
three probabilities, and the joint law is trinomial. That the generating function
factorises in this way is \cref{m:eq:indRelPGF} read backwards.

The absorption--birth--death generating function of \cref{m:prop:Gabd} is not of the
form \cref{m:eq:multilinear}, and there \cref{m:eq:stateProb} must be applied to
\cref{m:eq:Gabd} directly, or the coefficients extracted numerically by Cauchy
integrals on circles $|x|=|y|=\varrho<1$.

% --- restore default section numbering for subsequent chapters --------------
\renewcommand{\thesection}{\thechapter.\arabic{section}}
